\documentclass[10pt,a4paper]{amsart}
\usepackage{amsmath,amssymb,amsthm,graphicx}
% Retained preamble for the Stacks Project chapter snapshot in this directory.
% The chapter's own preamble.tex is a copy of the shared Stacks Project preamble; the block
% below is sliced MECHANICALLY, line for line, from it. Nothing was added, removed or
% redefined; the wrapper class, the xr-hyper/hyperref wiring and the externaldocument list
% are omitted because the assembler supplies the document class.
%
% lines taken from preamble.tex: 3, 5, 9, 11, 12, 156, 157, 158, ... (51 lines total)
%
% package coverage restored from the chapter's own preamble (2026-10-13): xy, xy-2cell

\usepackage[all]{xy}
\xyoption{2cell}
\UseAllTwocells

\newenvironment{reference}{\comment}{\endcomment}
\newenvironment{slogan}{\comment}{\endcomment}
\newenvironment{history}{\comment}{\endcomment}
\theoremstyle{plain}
\newtheorem{theorem}[subsection]{Theorem}
\newtheorem{proposition}[subsection]{Proposition}
\newtheorem{lemma}[subsection]{Lemma}
\theoremstyle{definition}
\newtheorem{definition}[subsection]{Definition}
\newtheorem{example}[subsection]{Example}
\newtheorem{exercise}[subsection]{Exercise}
\newtheorem{situation}[subsection]{Situation}
\theoremstyle{remark}
\newtheorem{remark}[subsection]{Remark}
\newtheorem{remarks}[subsection]{Remarks}
\numberwithin{equation}{subsection}
\def\lim{\mathop{\mathrm{lim}}\nolimits}
\def\colim{\mathop{\mathrm{colim}}\nolimits}
\def\Spec{\mathop{\mathrm{Spec}}}
\def\Hom{\mathop{\mathrm{Hom}}\nolimits}
\def\Ext{\mathop{\mathrm{Ext}}\nolimits}
\def\SheafHom{\mathop{\mathcal{H}\!\mathit{om}}\nolimits}
\def\SheafExt{\mathop{\mathcal{E}\!\mathit{xt}}\nolimits}
\def\Sch{\mathit{Sch}}
\def\Mor{\mathop{\mathrm{Mor}}\nolimits}
\def\Ob{\mathop{\mathrm{Ob}}\nolimits}
\def\Sh{\mathop{\mathit{Sh}}\nolimits}
\def\NL{\mathop{N\!L}\nolimits}
\def\CH{\mathop{\mathrm{CH}}\nolimits}
\def\proetale{{pro\text{-}\acute{e}tale}}
\def\etale{{\acute{e}tale}}
\def\QCoh{\mathit{QCoh}}
\def\Ker{\mathop{\mathrm{Ker}}}
\def\Im{\mathop{\mathrm{Im}}}
\def\Coker{\mathop{\mathrm{Coker}}}
\def\Coim{\mathop{\mathrm{Coim}}}
\DeclareMathSymbol{\boxtimes}{\mathbin}{AMSa}{"02}
\def\QCohstack{\mathcal{QC}\!\mathit{oh}}
\def\Cohstack{\mathcal{C}\!\mathit{oh}}
\def\Spacesstack{\mathcal{S}\!\mathit{paces}}
\def\Quotfunctor{\mathrm{Quot}}
\def\Hilbfunctor{\mathrm{Hilb}}
\def\Curvesstack{\mathcal{C}\!\mathit{urves}}
\def\Polarizedstack{\mathcal{P}\!\mathit{olarized}}
\def\Complexesstack{\mathcal{C}\!\mathit{omplexes}}
\def\Pic{\mathop{\mathrm{Pic}}\nolimits}
\def\Picardstack{\mathcal{P}\!\mathit{ic}}
\def\Picardfunctor{\mathrm{Pic}}
\def\Deformationcategory{\mathcal{D}\!\mathit{ef}}
% lane-s11: discard bibliographic \begin{reference} blocks entirely (incl. any \cite inside).
\newsavebox{\limrefbox}
\ifcsname reference\endcsname\expandafter\let\csname reference\endcsname\undefined\fi\ifcsname endreference\endcsname\expandafter\let\csname endreference\endcsname\undefined\fi
\def\citation#1{\relax}
\newenvironment{reference}{\setbox\limrefbox=\vbox\bgroup\let\cite\relax}{\egroup\ignorespaces}
\usepackage{cleveref}
\crefname{originalreference}{original reference}{original references}
\Crefname{originalreference}{Original reference}{Original references}
\makeatletter
\newlabel{algebra-lemma-characterize-finite-presentation}{{00QO}{source}{}{sourcepointer-1}{}}
\newlabel{algebra-lemma-characterize-finite-presentation@cref}{{[originalreference][0][]00QO}{[1][0][]source}{}{}{}}
\makeatother
\crefname{originallemma}{original lemma}{original lemmas}
\Crefname{originallemma}{Original lemma}{Original lemmas}
\makeatletter
\newlabel{lemma-limit-closed-nonempty}{{01Z3}{source}{}{sourcepointer-2}{}}
\newlabel{lemma-limit-closed-nonempty@cref}{{[originallemma][0][]01Z3}{[1][0][]source}{}{}{}}
\makeatother
\crefname{originalreference}{original reference}{original references}
\Crefname{originalreference}{Original reference}{Original references}
\makeatletter
\newlabel{schemes-lemma-diagonal-immersion}{{01KJ}{source}{}{sourcepointer-3}{}}
\newlabel{schemes-lemma-diagonal-immersion@cref}{{[originalreference][0][]01KJ}{[1][0][]source}{}{}{}}
\makeatother
\crefname{originalreference}{original reference}{original references}
\Crefname{originalreference}{Original reference}{Original references}
\makeatletter
\newlabel{morphisms-lemma-locally-finite-presentation-characterize}{{01TQ}{source}{}{sourcepointer-4}{}}
\newlabel{morphisms-lemma-locally-finite-presentation-characterize@cref}{{[originalreference][0][]01TQ}{[1][0][]source}{}{}{}}
\makeatother
\crefname{originallemma}{original lemma}{original lemmas}
\Crefname{originallemma}{Original lemma}{Original lemmas}
\makeatletter
\newlabel{lemma-descend-opens}{{01Z4}{source}{}{sourcepointer-5}{}}
\newlabel{lemma-descend-opens@cref}{{[originallemma][0][]01Z4}{[1][0][]source}{}{}{}}
\makeatother
\crefname{originallemma}{original lemma}{original lemmas}
\Crefname{originallemma}{Original lemma}{Original lemmas}
\makeatletter
\newlabel{lemma-limit-affine}{{01Z6}{source}{}{sourcepointer-6}{}}
\newlabel{lemma-limit-affine@cref}{{[originallemma][0][]01Z6}{[1][0][]source}{}{}{}}
\makeatother
\title{Three characterizations of locally finite presentation via limits of affine schemes}
\author{The Stacks Project Authors}
\date{stacks-project:limits.tex:local-snapshot-20261012; GNU FDL 1.2 or later}
\begin{document}
\maketitle
\noindent\textit{Editorial note.} EGA IV Proposition 8.14.2 cited in the source reference block.\par\medskip
\noindent\textit{Editorial note.} No omissions, hints or 'immediate' steps: all three implications are fully worked in the counted span.\par\medskip
\noindent\textit{Allowed standard prerequisites (editorial).}\begin{itemize}
\item Algebraic characterization of finite presentation. Algebra, Lemma algebra-lemma-characterize-finite-presentation
\item Limit closed nonempty principle. Limits, Lemma lemma-limit-closed-nonempty (same chapter)
\end{itemize}\medskip
\ifdefined\hypertarget\hypertarget{sourcepointer-1}{}\fi
\noindent\textit{Source reference.} Original reference 00QO: sources/stacks-chow/tags.txt, source lines 974--974. Cross-chapter standard prerequisite cited by the author inside the counted span; it is used as a pointer and its own proof is not part of this item.\par\medskip
\ifdefined\hypertarget\hypertarget{sourcepointer-2}{}\fi
\noindent\textit{Source reference.} Original lemma 01Z3: sources/stacks-chow/tags.txt, source lines 2578--2578. Cross-chapter standard prerequisite cited by the author inside the counted span; it is used as a pointer and its own proof is not part of this item.\par\medskip
\ifdefined\hypertarget\hypertarget{sourcepointer-3}{}\fi
\noindent\textit{Source reference.} Original reference 01KJ: sources/stacks-chow/tags.txt, source lines 2052--2052. Cross-chapter standard prerequisite cited by the author inside the counted span; it is used as a pointer and its own proof is not part of this item.\par\medskip
\ifdefined\hypertarget\hypertarget{sourcepointer-4}{}\fi
\noindent\textit{Source reference.} Original reference 01TQ: sources/stacks-chow/tags.txt, source lines 2383--2383. Cross-chapter standard prerequisite cited by the author inside the counted span; it is used as a pointer and its own proof is not part of this item.\par\medskip
\ifdefined\hypertarget\hypertarget{sourcepointer-5}{}\fi
\noindent\textit{Source reference.} Original lemma 01Z4: sources/stacks-chow/tags.txt, source lines 2579--2579. Cross-chapter standard prerequisite cited by the author inside the counted span; it is used as a pointer and its own proof is not part of this item.\par\medskip
\ifdefined\hypertarget\hypertarget{sourcepointer-6}{}\fi
\noindent\textit{Source reference.} Original lemma 01Z6: sources/stacks-chow/tags.txt, source lines 2581--2581. Cross-chapter standard prerequisite cited by the author inside the counted span; it is used as a pointer and its own proof is not part of this item.\par\medskip
\begin{proposition}
\label{proposition-characterize-locally-finite-presentation}
\begin{reference}
\cite[IV, Proposition 8.14.2]{EGA}
\end{reference}
Let $f : X \to S$ be a morphism of schemes.
The following are equivalent:
\begin{enumerate}
\item The morphism $f$ is locally of finite presentation.
\item For any directed set $I$, and any
inverse system $(T_i, f_{ii'})$ of $S$-schemes over $I$
with each $T_i$ affine, we have
$$
\Mor_S(\lim_i T_i, X) =
\colim_i \Mor_S(T_i, X)
$$
\item For any directed set $I$, and any
inverse system $(T_i, f_{ii'})$ of $S$-schemes over $I$
with each $f_{ii'}$ affine and every $T_i$ quasi-compact and
quasi-separated as a scheme, we have
$$
\Mor_S(\lim_i T_i, X) =
\colim_i \Mor_S(T_i, X)
$$
\end{enumerate}
\end{proposition}

\begin{proof}
It is clear that (3) implies (2).

\medskip\noindent
Let us prove that (2) implies (1). Assume (2).
Choose any affine opens $U \subset X$ and $V \subset S$ such that
$f(U) \subset V$. We have to show that
$\mathcal{O}_S(V) \to \mathcal{O}_X(U)$ is of finite presentation.
Let $(A_i, \varphi_{ii'})$ be a directed system of
$\mathcal{O}_S(V)$-algebras. Set $A = \colim_i A_i$.
According to
Algebra, Lemma \ref{algebra-lemma-characterize-finite-presentation}
we have to show that
$$
\Hom_{\mathcal{O}_S(V)}(\mathcal{O}_X(U), A) =
\colim_i \Hom_{\mathcal{O}_S(V)}(\mathcal{O}_X(U), A_i)
$$
Consider the schemes $T_i = \Spec(A_i)$. They
form an inverse system of $V$-schemes over $I$
with transition morphisms $f_{ii'} : T_i \to T_{i'}$
induced by the $\mathcal{O}_S(V)$-algebra maps $\varphi_{i'i}$.
Set $T := \Spec(A) = \lim_i T_i$.
The formula above becomes in terms of morphism sets of schemes
$$
\Mor_V(\lim_i T_i, U) =
\colim_i \Mor_V(T_i, U).
$$
We first observe that
$\Mor_V(T_i, U) = \Mor_S(T_i, U)$
and
$\Mor_V(T, U) = \Mor_S(T, U)$.
Hence we have to show that
$$
\Mor_S(\lim_i T_i, U) =
\colim_i \Mor_S(T_i, U)
$$
and we are given that
$$
\Mor_S(\lim_i T_i, X) =
\colim_i \Mor_S(T_i, X).
$$
Hence it suffices to prove that given a morphism $g_i : T_i \to X$ over $S$
such that the composition $T \to T_i \to X$ ends up in $U$ there exists some
$i' \geq i$ such that the composition $g_{i'} : T_{i'} \to T_i \to X$ ends up
in $U$. Denote $Z_{i'} = g_{i'}^{-1}(X \setminus U)$.
Assume each $Z_{i'}$ is nonempty
to get a contradiction. By Lemma \ref{lemma-limit-closed-nonempty}
there exists a point $t$ of $T$ which is mapped into $Z_{i'}$ for all
$i' \geq i$. Such a point is not mapped into $U$. A contradiction.

\medskip\noindent
Finally, let us prove that (1) implies (3). Assume (1). Let an inverse directed
system $(T_i, f_{ii'})$ of $S$-schemes be given. Assume the morphisms $f_{ii'}$
are affine and each $T_i$ is quasi-compact and quasi-separated as a scheme. Let
$T = \lim_i T_i$. Denote $f_i : T \to T_i$ the projection morphisms.
We have to show:
\begin{enumerate}
\item[(a)] Given morphisms $g_i, g'_i : T_i \to X$ over $S$ such that
$g_i \circ f_i = g'_i \circ f_i$, then there exists an $i' \geq i$
such that $g_i \circ f_{i'i} = g'_i \circ f_{i'i}$.
\item[(b)] Given any morphism
$g : T \to X$ over $S$ there exists an $i \in I$ and a morphism
$g_i : T_i \to X$ such that $g = f_i \circ g_i$.
\end{enumerate}

\noindent
First let us prove the uniqueness part (a). Let $g_i, g'_i : T_i \to X$ be
morphisms such that $g_i \circ f_i = g'_i \circ f_i$. For any $i' \geq i$
we set $g_{i'} = g_i \circ f_{i'i}$ and $g'_{i'} = g'_i \circ f_{i'i}$.
We also set $g = g_i \circ f_i = g'_i \circ f_i$.
Consider the morphism
$(g_i, g'_i) : T_i \to X \times_S X$. Set
$$
W =
\bigcup\nolimits_{U \subset X\text{ affine open},
V \subset S\text{ affine open}, f(U) \subset V}
U \times_V U.
$$
This is an open in $X \times_S X$, with the property that the morphism
$\Delta_{X/S}$ factors through a closed immersion into $W$, see the proof
of Schemes, Lemma \ref{schemes-lemma-diagonal-immersion}.
Note that the composition
$(g_i, g'_i) \circ f_i : T \to X \times_S X$ is a morphism into $W$
because it factors through the diagonal by assumption.
Set $Z_{i'} = (g_{i'}, g'_{i'})^{-1}(X \times_S X \setminus W)$.
If each $Z_{i'}$ is nonempty, then by Lemma \ref{lemma-limit-closed-nonempty}
there exists a point $t \in T$ which maps to $Z_{i'}$ for all
$i' \geq i$. This is a contradiction with the fact that $T$ maps into $W$.
Hence we may increase $i$ and assume that $(g_i, g'_i) : T_i \to X \times_S X$
is a morphism into $W$. By construction of $W$, and since $T_i$ is
quasi-compact we can find a finite affine open covering
$T_i = T_{1, i} \cup \ldots \cup T_{n, i}$ such that
$(g_i, g'_i)|_{T_{j, i}}$ is a morphism into $U \times_V U$ for
some pair $(U, V)$ as in the definition of $W$ above.
Since it suffices to prove that $g_{i'}$ and $g'_{i'}$ agree
on each of the $f_{i'i}^{-1}(T_{j, i})$ this reduces us to the affine case.
The affine case follows from
Algebra, Lemma \ref{algebra-lemma-characterize-finite-presentation}
and the fact that the ring map
$\mathcal{O}_S(V) \to \mathcal{O}_X(U)$ is of finite presentation
(see Morphisms,
Lemma \ref{morphisms-lemma-locally-finite-presentation-characterize}).

\medskip\noindent
Finally, we prove the existence part (b).
Let $g : T \to X$ be a morphism of schemes over $S$.
We can find a finite affine open covering
$T = W_1 \cup \ldots \cup W_n$ such that for
each $j \in \{1, \ldots, n\}$ there exist affine opens
$U_j \subset X$ and $V_j \subset S$ with $f(U_j) \subset V_j$
and $g(W_j) \subset U_j$. By Lemmas \ref{lemma-descend-opens}
and \ref{lemma-limit-affine}
(after possibly shrinking $I$) we may assume that
there exist affine open coverings $T_i = W_{1, i} \cup \ldots \cup W_{n, i}$
compatible with transition maps such that $W_j = \lim_i W_{j, i}$.
We apply Algebra, Lemma \ref{algebra-lemma-characterize-finite-presentation}
to the rings corresponding to the affine schemes $U_j$, $V_j$, $W_{j, i}$ and
$W_j$ using that $\mathcal{O}_S(V_j) \to \mathcal{O}_X(U_j)$ is of finite
presentation (see Morphisms,
Lemma \ref{morphisms-lemma-locally-finite-presentation-characterize}).
Thus we can find for each $j$ an index $i_j \in I$ and a morphism
$g_{j, i_j} : W_{j, i_j} \to X$ such that
$g_{j, i_j} \circ f_i|_{W_j} : W_j \to W_{j, i} \to X$
equals $g|_{W_j}$. By part (a) proved above, using the quasi-compactness of
$W_{j_1, i} \cap W_{j_2, i}$ which follows as $T_i$ is quasi-separated,
we can find an index $i' \in I$ larger than all $i_j$ such that
$$
g_{j_1, i_{j_1}} \circ f_{i'i_{j_1}}|_{W_{j_1, i'} \cap W_{j_2, i'}} =
g_{j_2, i_{j_2}} \circ f_{i'i_{j_2}}|_{W_{j_1, i'} \cap W_{j_2, i'}}
$$
for all $j_1, j_2 \in \{1, \ldots, n\}$. Hence the morphisms
$g_{j, i_j} \circ f_{i'i_j}|_{W_{j, i'}}$ glue to given the
desired morphism $T_{i'} \to X$.
\end{proof}

\end{document}
